%!TEX root = ./../../springer-mono-format/book.tex
%!TEX root = ./../../thesis-trier.tex

Permutation group theory enters graph theory via the automorphism group 
in a natural way. Here, we introduce some further notions
connecting graphs and permutation groups. \Cref{prelim:sec:genSetSn} collects results on generating sets for $\mathcal S_n$. For example, we give a graph-theoretical
characterization when a set of transpositions generate~$\mathcal S_n$.
\Cref{prelim:sec:vertex-transitive}
introduces vertex-transitive graphs and investigates the connectivity
of vertex-transitive graphs. 
%\Cref{prelim:thm:degree-implies-connectivity} \todo{finite?}  
%links the outdegree of a strongly connected vertex-transitive
%graph to its strong $k$-connectivity. 
Lastly, \cref{prelim:sec:orbitalgraphs}
introduces orbital graphs.

We often denote graphs and groups
with the capital letter~$G$, but the context
should make it clear whether we mean a graph or a group.


\subsection{Generating Sets for the Full Symmetric Group}
\label{prelim:sec:genSetSn}

It is a well-known result that permutation cyclically permuting a finite set
and a transposition on this set generate the full symmetric group.
Isaacs and Zieschang~\cite[Thm.~B]{IsaacsZ95} gave the following generalization.

\begin{theorem}[{Isaacs and Zieschang~\cite[Thm.~B]{IsaacsZ95}}]
	\label{prelim:thm:cycles_gen_Sn}
	Let $x = (0 ~ 1 ~ 2 ~ \cdots ~ n-1)_{[n]}$ and $y = (0 ~ 1 ~ 2 ~ \cdots ~ m-1)_{[n]}$ for some $1 < m < n$.
	Then $\langle \{x,y\} \rangle_G$ equals the full symmetric group on $[n]$ unless both $n$ and $m$ are odd, in which
	case $\langle \{x,y\} \rangle_G$ equals $\mathcal A_n$. % is the alternating group.
\end{theorem}

Isaacs and Zieschang~\cite[Thm.~A]{IsaacsZ95} also proved that for $n \ne 4$
and $x \in \mathcal S_n$ there exists $y \in \mathcal S_n$ such that $\langle \{x,y\} \rangle_G = \mathcal S_n$  (the exception being any element $x$ in the Klein four-group from \cref{prelim:permgrp:ex:Kleinfourgroup}).
Actually, this result was first proved by Sophie Piccard~\cite[pp.~80--86]{piccard46} in 1946, 
but got lost in limbo\footnote{We were not able to obtain
	a copy of Piccards original work~\cite{piccard46}
	but learned this through the work of Salomaa~\cite{DBLP:journals/jucs/Salomaa02}.}
and was reproved by Isaacs \& Zieschang~\cite{IsaacsZ95} in 1996.

By \cref{prelim:fact:transpositions_generate_Sn} the full symmetric group on a finite set
is generated by all transpositions. Next, we state a refined result, needed in \cref{sync_max:sec:perm_grps},
that characterizes the sets of transpositions generating $\mathcal S_n$ in terms of an associated graph. 
For the convenience of the reader, we give a proof (which can also be found in~\cite[Lem.~3.10.1]{godsil01};
compare also the proof of~\cite[Lem.~6(a)]{AraujoCameron2014}
which uses a similar argument).

\begin{proposition}
	\label{prop:connected_transpositions_gen_Sn}
	Let $T \subseteq \mathcal S_n$ be a set of transpositions
	and consider the undirected graph $G = ([n], \{ \{a,b\} \mid (a ~ b) \in T\})$.
	Then $G$ is connected if and only if $\langle T \rangle_G = \mathcal S_n$.
\end{proposition}
\begin{proof}
	Define a relation $\sim$ on $[n]$
	by $a \sim b$ if and only if $a = b$ or $(a ~ b) \in \langle T \rangle_G$.
	Clearly, $\sim$ is reflexive and symmetric.
	Further, for $a \sim b$ and $b \sim c$
	with distinct $a,b,c \in [n]$ it is $(a ~ b)^{(b ~ c)} = (a ~ c) \in \langle T \rangle_G$;
	thus $\sim$ is transitive and an equivalence relation.
	As the graph $G$ results from the digraph $G_{\sim}$ by forgetting
	about the direction of the edges, from \cref{prelim:fact:relation_graph}
	we obtain that the connected components equal the equivalence classes of~$\sim$.	
	So, if $G$ is connected, then all points are pairwise comparable for $\sim$, and
	$\langle T \rangle_G$ contains all transpositions.
	By Fact~\ref{prelim:fact:transpositions_generate_Sn}, we find $\langle T \rangle_G = \mathcal S_n$.
	Otherwise, suppose~$G$ is not connected.
	Then, we can pick $a,b$
	with $a \not\sim b$, i.e., $(a ~ b) \notin \langle T \rangle_G$
	and so~$\langle T \rangle_G \ne \mathcal S_n$.
\end{proof}

\begin{remark}	
	In the proof of Jordan's theorem from~\cite{IsaacsZ95},
	it is also shown that for a set of $3$-cycles $T$, $\langle T \rangle_G$ contains the alternating
	group if and only if the graph on~$[n]$
	where $a,b \in [n]$ are joined by an edge if and only if 
	$a,b \in \move{t}$ for some $t \in T$
	is connected.
\end{remark}

In~\cite{Knoerr2007} it was shown that a $k$-transitive group of degree $n > k$
is generated by all elements fixing precisely~$i$ points if $i \le k - 3$.
Since the full symmetric group on a set of size~$n$ is $k$-transitive for each $1 \le k \le n$,
by this result we obtain more generating sets for the full symmetric group of a finite set.

\subsection{Vertex-Transitive Graphs}
\label{prelim:sec:vertex-transitive}

\IN{vertex-transitive graph}
\IN{outdegree of graph}
\IN{strong connectivity $\kappa(G)$}
\label{def:kappa}
A (directed or undirected) graph is \emph{vertex-transitive} if its automorphism group acts transitively on its
vertices.
If~$G = (V,E)$ is vertex-transitive, then every vertex has the same outdegree 
and we call the common value the \emph{outdegree of~$G$}.
%The \emph{strong connectivity}~$\kappa(G)$ is the unique number~$k$ \todo{warum $K_n$ ausgeklammert? weil dann nie zwei SCCs}
%such that % K_n sagen einfach teil der ausagen.
%$G$ is strongly~$k$-connected but not strongly $(k+1)$-connected.

A proof of the next statement can be found in~\cite[Thm~1.10]{cameron_1999},
\cite[Thm~1.1]{ArCamSt2017}, \cite[Chap.~2.6]{godsil01},
and in~\cite{Rystsov2022} (in the latter work formulated with slightly different notions but easily seen to give the result).

\begin{theorem}
	\label{prelim:thm:vertex_transitive_finite_graph}
	Let $G = (V, E)$ be a finite vertex-transitive digraph.
	If $G$ is weakly connected, then $G$ is strongly connected.
\end{theorem}

Next, we link the outdegree of a vertex-transitive digraph to its connectivity. The result appears in~\cite[Prop.~3.4]{Hamidoune81} (compare also~\cite[Thm~3.7 (a)]{BabaiSurvey96} where it is formulated
in terms of the connectivity 
from \cref{prelim:rem:connectedconnectivity}).

\begin{theorem}[{Hamidoune~\cite{Hamidoune81}}]
	\label{prelim:thm:degree-implies-connectivity}
	If $G = (V, E)$
	is a finite strongly connected vertex-transitive digraph of outdegree~$d$ without edges of the form $(v, v)$ for $v \in V$,
	then~$G$ is strongly
	$\left( \left\lfloor \frac{d}{2} \right\rfloor + 1 \right)$-connected.
\end{theorem} 
\begin{proof} \label{def:Nplus} \IN{atom in a graph}
	We can assume $E \ne V \times V$. 
	Since~$G$ is finite, there exists a largest number~$k$ 
	such that~$G$ is 
	strongly $k$-connected, but not strongly $(k+1)$-connected.
	For $U \subseteq V$
	set $N^+(U) = \{ v \in V \setminus U \mid \exists u \in U : (u,v) \in E \}$. 
	We have $k \le |N^+(\{u\})| = d$ for each $u \in V$. % A = N^+(u) muss jeder pfad zu nem anderen schneiden, also nicht ( |N^+(u)|+1 )-connected, damit | | \le kappe
	An \emph{atom}~\cite{hamidoune77} $A \subseteq V$ is a set of minimal cardinality with
	$|N^+(A)| = k$ and
	$A \cup N^+(A) \ne V$.\footnote{For undirected graphs, atoms 
		were introduced in~\cite{Mader70,Mader71,Watkins70}.}
	In a strongly connected digraph the atoms partition~$V$~\cite{hamidoune77}
	(see also~\cite[Thm. 2.1 in both]{DBLP:journals/ejc/Hamidoune81,Hamidoune81}).
	As an automorphism maps atoms to atoms,
	the atoms induce an $\operatorname{Aut}(G)$-congruence on the vertices.
	Hamidoune proved that $|A| \le k$~\cite[Lem.~3.3]{Hamidoune81}.
	So, for $u \in A$, as $N^+(\{u\}) \subseteq A \setminus \{u\} \cup N^+(A)$, we get 
	%	$ 
	%	d \le |A \setminus \{u\}| + |N^+(\{u\})| < |A| + |N^+(A)| = |A| + \kappa \le 2\kappa. 
	%	$
	% N^+ nicht monoton bzgl submengen (beispiel dreieck mit allen richtungen und 2-er menge und 3-er menge)
	$
	d \le |A \setminus \{u\}| + |N^+(A)| < |A| + k \le 2k.
	$ % u betrachten, geht entweder raus, oder bleibt in A zu den A \ {u} anderen, damit folgt erste Ungleichung.
\end{proof} 

In the context of the present work, we actually only need the following
implication.

\begin{corollary}
	\label{prelim:prop:vertex-transitive-SCC-32-connected}
	If $G = (V, E)$
	is a finite strongly connected vertex-transitive graph without edges of the form $(v,v)$ for $v \in V$,
	then $G$ is strongly $\frac{3}{2}$-connected.
\end{corollary}
\begin{proof}
	Let $d$ be the outdegree of $G$.	
	For $d = 1$, $G$ forms a single cycle,
	i.e., a strongly connected digraph
	where every vertex has outdegree one,	
	which is clearly strongly~$\frac{3}{2}$-connected.
	For $d \ge 2$,
	by \cref{prelim:thm:degree-implies-connectivity}, 
	$G$ is strongly~$2$-connected, hence strongly $\frac{3}{2}$-connected.	
\end{proof} 

Note that, by \cref{prelim:prop:vertex-transitive-SCC-32-connected},
the conclusion of \cref{prelim:thm:degree-implies-connectivity}
is valid for $d = 1$ even after removal of the 
rounding, which somehow explains
our choice of naming for strong $\frac{3}{2}$-connectivity.

\subsection{Orbital Graphs}
\label{prelim:sec:orbitalgraphs}

Suppose $G \le_G \mathcal S_X$ is a permutation group.
Let $a,b \in X$. Then we can associate two graphs with the tuple $(a,b)$
in a natural way.

\IN{orbital graph!directed}
\IN{orbital graph!undirected}
\begin{enumerate} 
	\item The \emph{directed orbital graph} of the tuple $(a,b)$
	is the graph with vertex set~$X$ and edge set $\{ (\applyfunc{a}{g}, \applyfunc{b}{g}) \mid g \in G \}$.
	
	\item The \emph{undirected orbital graph} of the tuple $(a,b)$ (or the $2$-set $\{a,b\}$)
	is the graph with vertex set~$X$ and edge set $\{ \{ \applyfunc{a}{g}, \applyfunc{b}{g} \} \mid g \in G \}$.	
\end{enumerate} 

\begin{example} 
 \Cref{fig:orbitals} shows the directed orbital graphs
 of the group $G = \langle \{ (0 ~ 1 ~ 2 ~ 3)\} \rangle_G$.
 The undirected orbital graphs are obtained from these
 by forgetting about the directions on the edges.
 Since~$G$ contains no element swapping two points, 
 the orbital graphs are oriented.

\todo[inline]{graphiken zu groß}

\begin{figure}[th]
	\centering
	\begin{tikzpicture}[>=Stealth, node distance=2em, every node/.style={minimum size=0.5cm}]
		\node (1) at (0,2) {0};
		\node (2) at (2,2) {1};
		\node (3) at (2,0) {2};
		\node (4) at (0,0) {3};
		
		\path[->] (1) edge [loop above] (1);
		\path[->] (2) edge [loop above] (2);
		\path[->] (3) edge [loop above] (3);
		\path[->] (4) edge [loop above] (4);
		
		\node (1r) at (4,2) {0};
		\node (2r) at (6,2) {1};
		\node (3r) at (6,0) {2};
		\node (4r) at (4,0) {3};
		
		\draw (1r) -- (2r); \draw (2r) -- (3r); \draw (3r) -- (4r); \draw (4r) -- (1r);
				
		\node (1s) at (8,2) {0};
		\node (2s) at (10,2) {1};
		\node (3s) at (10,0) {2};
		\node (4s) at (8,0) {3};
		
		\draw (1s) -- (3s); \draw (4s) -- (2s);
		\draw (3s) -- (1s); \draw (2s) -- (4s);
		
		\node (1t) at (12,2) {0};
		\node (2t) at (14,2) {1};
		\node (3t) at (14,0) {2};
		\node (4t) at (12,0) {3};
		
		\draw (1t) -- (4t); \draw (4t) -- (3t); \draw (3t) -- (2t); \draw (2t) -- (1t);
		
	\end{tikzpicture}
	\caption[Directed orbital graphs of $G = \langle \{ (0 ~ 1 ~ 2 ~ 3)\} \rangle_G$]{The directed orbital graphs of $G = \langle \{ (0 ~ 1 ~ 2 ~ 3)\} \rangle_G$.}
	\label{fig:orbitals}
\end{figure}
\end{example}

\begin{example}
	\Cref{fig:orbitals-klein-four} shows the directed orbital graphs
	of the group $G = \langle \{ (0 ~ 1)(2 ~ 3), (0 ~ 2)(1 ~ 3)\} \rangle_G$.
	All of them except the left-most one are symmetric.
	
	\begin{figure}[th]
		\centering
		\begin{tikzpicture}[>=triangle 45, node distance=1.5cm, every node/.style={minimum size=0.5cm}]
			
			\node (1) at (0,2) {0};
			\node (2) at (2,2) {1};
			\node (3) at (2,0) {2};
			\node (4) at (0,0) {3};
			
			\path[->] (1) edge [loop above] (1);
			\path[->] (2) edge [loop above] (2);
			\path[->] (3) edge [loop above] (3);
			\path[->] (4) edge [loop above] (4);
			
			\node (1r) at (4,2) {0};
			\node (2r) at (6,2) {1};
			\node (3r) at (6,0) {2};
			\node (4r) at (4,0) {3};
			
			\draw[<->] (1r) -- (2r);
			\draw[<->] (3r) -- (4r);
			
			\node (1s) at (8,2) {0};
			\node (2s) at (10,2) {1};
			\node (3s) at (10,0) {2};
			\node (4s) at (8,0) {3};
			
			\draw[<->] (1s) -- (3s);
			\draw[<->] (2s) -- (4s);
						
			\node (1t) at (12,2) {0};
			\node (2t) at (14,2) {1};
			\node (3t) at (14,0) {2};
			\node (4t) at (12,0) {3};
			
			\draw[<->] (1t) -- (4t);
			\draw[<->] (2t) -- (3t);
			
		\end{tikzpicture}
		\caption[Directed orbital graphs of $G = \langle \{ (0 ~ 1)(2 ~ 3), (0 ~ 2)(1 ~ 3)\} \rangle_G$]{The directed orbital graphs of $G = \langle \{ (0 ~ 1)(2 ~ 3), (0 ~ 2)(1 ~ 3)\} \rangle_G$.}
		\label{fig:orbitals-klein-four}
	\end{figure}
\end{example}

The edge set of the directed orbital graph of $(a,b)$
is contained in $\Orbit{a}{G} \times \Orbit{b}{G}$.
The directed orbital graph of $(a,b)$ 
is either symmetric
or oriented, and the former case applies precisely if $G$
contains a permutation swapping~$a$ and~$b$~\cite[Lem.~2.4.2]{godsil01}.
Also, $G$ is intransitive if and only if 
in the (directed or undirected) orbital graph
exist vertices that are not part of any edge.
Furthermore,~$G$ is a subgroup of the automorphism
group of each (directed or undirected) orbital graph.
So, by Theorem~\ref{prelim:thm:vertex_transitive_finite_graph} 
the weakly connected components of a finite directed orbital graph are precisely
the strongly connected components.
Finally, a (directed or undirected)
graph is vertex-transitive
if and only if it is a union of orbital graphs w.r.t.~its automorphism group.

\begin{lemma} \label{prelim:lem:orbitals-distinct-orbits}
	Let $G \le_G \mathcal S_n$ and $a,b \in [n]$.
	Then $\Orbit{a}{G} \ne \Orbit{b}{G}$
	if and only if in the undirected orbital graph of $\{a,b\}$ (resp.~directed orbital graph of $(a,b)$)
	there exists no edge with 
	both 
	endpoints in $\Orbit{a}{G}$
	or both endpoints in $\Orbit{b}{G}$
	(resp.~the starting vertex of each edge is in $\Orbit{a}{G}$ 
	and its ending vertex is in $\Orbit{b}{G}$).
\end{lemma}
\begin{proof}
	Since $\Orbit{a}{G}$ and $\Orbit{b}{G}$ are orbits, we have $\Orbit{a}{G} = \Orbit{b}{G}$
	if and only if there exists $g \in G$ such that $g(a) \in \Orbit{b}{G}$ or $g(b) \in \Orbit{a}{G}$.
	So,  $\{g(a), g(b)\} \subseteq \Orbit{a}{G}$
	or $\{g(a), g(b)\} \subseteq \Orbit{b}{G}$ for some $g \in G$
	if and only if~$G$ is transitive, which
	gives the claim for the undirected orbital graph.
	Furthermore, since the edge set of the directed orbital graph is in $\Orbit{a}{G} \times \Orbit{b}{G}$,
	if $G$ is intransitive every edge of the directed orbital graph
	starts in $\Orbit{a}{G}$ and ends in $\Orbit{b}{G}$.	
\end{proof}